\section{The trace form and the residue data}
\label{trace:section}

The goal of this section is Proposition~\ref{prop:residue-data}: a
meromorphic pair with prescribed poles and two normalized residue
identities. The trace form translates the quadratic-space data on the
elliptic curve into those identities on $C$. The comparison must preserve
its scalar and sign, since both enter the first gluing obstruction.

Retain the data and the normalizations of Proposition~\ref{prop:curve-data}.
In particular, $\Omega$ is the fixed logarithmic volume on $\Sigma$, its
normal-first residue is the nowhere vanishing differential $\omega$ on
$\Gamma$, and $4L\simeq K_C$ is the specified adjunction isomorphism.
Put
\[
 \theta=2L,\qquad \Lambda=\theta+H=3H+2Z,
 \qquad f:C\longrightarrow\PP^1.
\]
For the chosen basis $u=(u_1,u_2)$ of the ruling pencil, write
$\eta_C=u_1\,du_2-u_2\,du_1\in H^0(C,K_C+2H)$.
Although the individual derivatives depend on a frame of $H$, this
alternating expression does not.  We use the analogous notation
$\eta_\Gamma$ on $\Gamma$.

\subsection{A two-dimensional quotient of a perfect trace form}

The pairing below is a concrete instance of finite-morphism duality,
as developed by Grothendieck and presented in Hartshorne's
\emph{Residues and Duality}~\cite[Chapter~III, \S6]{Hartshorne1966}. We check its local
perfection directly so that the trace normalization remains explicit.

\begin{lemma}\label{trace:pairing}
The rank-four vector bundle $E=f_*\Lambda$ has a perfect symmetric pairing
\begin{equation}\label{trace:form}
 E\otimes E\longrightarrow\OO_{\PP^1},\qquad
 (T_1,T_2)\longmapsto
 \Tr_{C/\PP^1}\!\left(\frac{T_1T_2}{\eta_C}\right).
\end{equation}
Its splitting type is $(3,0,0,-3)$.  Multiplication by $\xi^2$ identifies
$\OO_{\PP^1}(3)$ with its positive line subbundle $\mathcal P$.
The form on global sections has radical
$\xi^2H^0(C,3H)$, and induces a nondegenerate symmetric form
$\langle\ ,\ \rangle$ on
\begin{equation}\label{trace:W}
 W=H^0(C,\Lambda)/\xi^2H^0(C,3H),\qquad \dim W=2.
\end{equation}
For every $x\in\PP^1$, evaluation identifies this quadratic space with
$(\mathcal P^\perp/\mathcal P)_x$.

Suppose $q\in Z$, and choose local frames $\mathfrak f$ of $H$ and
$\tau$ of $\Lambda$, and a parameter $z$ at $q$.  Write, using the fixed
isomorphism $2\Lambda\simeq K_C+2H$,
\[
 \tau^2=s(z)\,dz\,\mathfrak f^2,
 \qquad \eta_C=\mathcal E(z)\,dz\,\mathfrak f^2.
\]
Coefficient evaluation $\ell_q:T\mapsto(T/\tau)(q)$ descends to $W$.
Its square in the dual form is
\begin{equation}\label{trace:dual-evaluation}
 \ell_q^2=\frac{\mathcal E(q)}{s(q)}.
\end{equation}
In particular, it is nonzero and nonisotropic.
\end{lemma}

\begin{proof}
The morphism $f$ is finite and flat of degree four.  The identity
$2\Lambda=K_C+2H$ identifies $\Lambda^2$ with the relative dualizing
bundle.  We verify directly the normalization and perfection of its trace
pairing.  Locally at a ramification point one can choose parameters with
$x=w^e$ on a single component over a target disk.  A frame of $H$ pulled
back from that disk can be chosen so that
$\eta_C=dx\,\mathfrak f^2$.  If
$\tau^2=a(w)\,dw\,\mathfrak f^2$, with $a$ a unit, the local pairing is
\[
 (b,c)\longmapsto
 \Tr\!\left(\frac{a(w)bc}{e w^{e-1}}\right).
\]
On the basis $1,w,\ldots,w^{e-1}$, the pairing without the unit $a$ has
matrix
\[
 \Tr\!\left(\frac{w^{i+j}}{e w^{e-1}}\right)
   =\begin{cases}1&i+j=e-1,\\0&i+j\ne e-1,
     \end{cases}
 \qquad 0\le i,j<e.
\]
It is therefore holomorphic and perfect.  Multiplication by $a$ is an
automorphism, and taking the direct sum over the points above the target
disk proves the assertion.  On an unramified fiber the trace is the sum
of its four contributions; there is no division by the degree.

By the splitting theorem for vector bundles on $\PP^1$
\cite[\S2, Theorem~2.1 and Corollary]{Grothendieck1957}, $E$ is a direct sum of line bundles.
Self-duality implies that its splitting degrees occur in opposite
pairs.  Proposition~\ref{prop:curve-data} gives
\[
 h^0(E(-1))=h^0(\theta)=3,\qquad
 h^0(E(-2))=h^0(H+2Z)=2.
\]
A positive splitting degree $a$ contributes $a$ and $a-1$, respectively,
to these two numbers, and all nonpositive degrees contribute zero.
There is consequently exactly one positive degree, equal to three.
Thus
\[
 E\simeq\OO(3)\oplus\OO^{\oplus2}\oplus\OO(-3).
\]
The nonzero map $\OO(3)\to E$ obtained from multiplication by $\xi^2$
has zero components in the three lower-degree summands.  Its component
in $\OO(3)$ is a nonzero constant.  Hence its image $\mathcal P$ is
the entire positive summand, and is a subbundle, including over the
images of the points of $Z$.
The injective pullback
$H^0(\PP^1,\OO(3))\to H^0(C,3H)$ is an isomorphism, since both
spaces have dimension four.  Thus
$H^0(\mathcal P)=\xi^2H^0(C,3H)$.

The line $\mathcal P$ is isotropic, since
$\Hom(\OO(6),\OO)=0$.  Likewise it pairs trivially with the two
degree-zero summands.  Perfection gives
\[
 \mathcal P^\perp\simeq\OO(3)\oplus\OO^{\oplus2},\qquad
 \mathcal P^\perp/\mathcal P\simeq\OO^{\oplus2},
\]
with a nondegenerate constant form on the quotient.  All global sections
of $E$ lie in $\mathcal P^\perp$, and $H^1(\mathcal P)=0$.
Taking sections proves the assertions about $W$ and its fibers.

At $q\in Z$, evaluation annihilates $\mathcal P_{f(q)}$, because the
map defining that line is multiplication by $\xi^2$.  The finite fiber algebra is a product of its Artin-local factors,
and the trace pairing is the orthogonal sum of their pairings.
Since $f$ is unramified at $q$, the factor at $q$ is one-dimensional,
with form coefficient $s(q)/\mathcal E(q)$. Evaluation at $q$
vanishes on the other factors, whether or not they are ramified.
The metric dual of coefficient evaluation thus has square
$\mathcal E(q)/s(q)$ in $E_{f(q)}$.  It lies in
$\mathcal P_{f(q)}^\perp$, because the evaluation kills
$\mathcal P_{f(q)}$.  Passing to the orthogonal quotient preserves this
square and gives the functional on $W$.  Finally
$\mathcal E(q)\ne0$ by the unramifiedness in
Proposition~\ref{prop:curve-data}.
\end{proof}

We write a dot for the dual bilinear form on $W^*$, so that
$\lambda^2=\lambda\cdot\lambda$.  The distinction between this dual
form and the form on $W$ accounts for the reciprocal in
\eqref{trace:dual-evaluation}.

\subsection{Identification by full principal parts}

The preceding lemma gives a quadratic space on $C$. We now identify it
with the one on $\Gamma$ used to choose the marked points. The
principal-part exact sequence, together with Serre duality
\cite[\S3, Theorem~4; \S4, Theorem~5]{SerreDuality1955}, characterizes
realizable principal parts by their residue pairings with global adjoint
sections.
We first compare both coefficients of each double pole to identify the
two spaces.  We then use the leading coefficients to determine the
scalar and sign of the quadratic form under that identification.

\begin{lemma}\label{trace:principal-parts}
There is an isomorphism
\begin{equation}\label{trace:iota}
 \iota:W\xrightarrow{\ \sim\ }H^0(\Gamma,\Phi)
\end{equation}
that identifies the full principal parts of $T/\xi^2$ on $C$ with
those of $\iota([T])/\beta^2$ on $\Gamma$ through residue duality.
It identifies the projective evaluations at corresponding points of $Z$.
Moreover, its dual
\[
 \iota^*:H^0(\Gamma,\Phi)^*\longrightarrow W^*
\]
is an isometry for the quadratic form $Q$ of
Proposition~\ref{prop:curve-data} and the dual trace form.
\end{lemma}

\begin{proof}
Set $N_\Sigma=\Gamma+F$, $N_\Gamma=N_\Sigma|_\Gamma$, and $D=2Z$.
Fourth-order contact identifies $D$ on $C$ and on $\Gamma$ as the same
embedded length-twelve subscheme of $\Sigma$.  The isomorphism
$N_\Sigma|_C\simeq K_C-3H$ is the one supplied by adjunction with the
fixed volume.  The differential $\omega$ trivializes $K_\Gamma$.
For a line bundle $B$ on a smooth curve, the local pairing between
principal parts of $B(D)$ and jets of $K-B$ is
\[
 (a,s)\longmapsto\sum_{q\in Z}\Res_q(as).
\]
It is perfect: in a parameter and mutually dual frames the two
coefficients of $a_{-2}z^{-2}+a_{-1}z^{-1}$ pair with a jet
$s_0+s_1z$ as $a_{-2}s_1+a_{-1}s_0$.
Consequently the two principal part spaces
\[
 H^0\!\left(C,\frac{\OO_C(3H+2Z)}{\OO_C(3H)}\right),
 \qquad
 H^0\!\left(\Gamma,
       \frac{\OO_\Gamma(-N_\Gamma+2Z)}{\OO_\Gamma(-N_\Gamma)}\right)
\]
are both identified by residues with
$H^0(D,N_\Sigma|_D)^*$.  This defines their comparison without any
choice of coordinates or frames. It compares the full local
principal-part spaces before any condition of global realizability
is imposed.

Both complete adjoint systems are restrictions of the same ambient
system $H^0(\Sigma,N_\Sigma)$.  Indeed, the kernels in the restriction
sequences to $C$ and to $\Gamma$ are respectively
\[
 \OO_\Sigma(-\Gamma-3F)=\OO_\Sigma(-5,-2),
 \qquad \OO_\Sigma(F)=\OO_\Sigma(1,0),
\]
and both have vanishing first cohomology.  Their images on the common
scheme $D$ therefore coincide.  Equivalently, an ambient section
vanishing on $\Gamma$ vanishes to order at least four on $C$ at $Z$
and has no effect on these residue tests.

The principal part exact sequence and Serre duality now say that the
image of $H^0(C,\Lambda)$, by $T\mapsto\operatorname{pp}(T/\xi^2)$,
is precisely the annihilator of this common adjoint jet space; its
kernel is $\xi^2H^0(C,3H)$.  On $\Gamma$ the same annihilator is the
image of
\[
 H^0(\Gamma,\Phi)\longrightarrow
 H^0\!\left(\Gamma,
       \frac{\OO_\Gamma(-N_\Gamma+2Z)}{\OO_\Gamma(-N_\Gamma)}\right),
 \qquad \widetilde T\longmapsto
       \operatorname{pp}(\widetilde T/\beta^2),
\]
because $\Phi=2M-N_\Gamma$ and $(\beta)=Z$.
This map is injective: its kernel would give a section of
$-N_\Gamma$, which has degree $-10$.
These facts establish \eqref{trace:iota}.  In particular, this is an
identification of full principal parts.  Its dimension can also be
read directly on $\Gamma$: $h^0(N_\Gamma)=10$, and restriction to
$D$ is injective since $\deg(N_\Gamma-2Z)=-2$; its annihilator has
dimension $12-10=2$.

We compute the local comparison, including its first derivative and the
adjunction scalar.  Near a point of $Z$ choose coordinates $(y,z)$ and
frames $\mathfrak g$ of $\Gamma$ and $\mathfrak f$ of $F$ such that
\[
 \gamma=y\mathfrak g,\qquad
 \Omega=e(y,z)\frac{dy}{y}\wedge dz,\qquad e(0,z)=1,
\]
so $\omega=dz$ on $\Gamma$.  Write
\[
 \kappa=k(y,z)\mathfrak g^2\mathfrak f^4,
 \qquad \beta=b(z)\mathfrak m,
 \qquad \mathfrak m^4=\mathfrak g^2\mathfrak f^4\big|_\Gamma.
\]
Here $k(0,z)=b(z)^4$ and
$b'(0)k_y(0,0)\ne0$.  With $K(z)=k_y(0,z)$, the equation of $C$ gives
\begin{equation}\label{trace:contact-expansion}
 y|_C=-\frac{b(z)^4}{K(z)}+O(z^8),\qquad
 e|_C=1+O(z^4),\qquad k_y|_C=K(z)+O(z^4).
\end{equation}
The normal-first adjunction convention sends the frame
$\mathfrak g\mathfrak f^4$ to
\begin{equation}\label{trace:adjunction}
 \Res_C\!\left(\frac{e(y,z)\,dy\wedge dz}{k(y,z)}\right)
       =\frac{e|_C}{k_y|_C}\,dz.
\end{equation}
The sign is positive here, since
$dy\wedge dz=(dk\wedge dz)/k_y$.

Let $a=\xi/z$ be the local frame of $\OO_C(Z)$ and
$\tau=\mathfrak f^3a^2$ the resulting frame of $\Lambda$.
The fixed fourth-power identification sends
$a^4$ to $(y|_C/z^4)\mathfrak g$.
Using \eqref{trace:adjunction}, its square has the form
\begin{equation}\label{trace:frame-square}
 \tau^2=s(z)\,dz\,\mathfrak f^2,
 \qquad
 s(z)=\frac{e|_C\,y|_C}{z^4 k_y|_C},\qquad
 s(0)=-\frac{b'(0)^4}{k_y(0,0)^2}.
\end{equation}
Thus the negative sign comes from solving the equation of $C$.

Write $T=\pi(z)\tau$, and write
$\widetilde T=\iota([T])$ in the frame
$\widetilde\tau=\mathfrak m^2/(\mathfrak g\mathfrak f)$ of $\Phi$
as $\widetilde T=\widetilde\pi(z)\widetilde\tau$.
For an ambient test section $A(y,z)\mathfrak g\mathfrak f$, the two
residue tests are
\begin{equation}\label{trace:local-residue-tests}
 \Res_0\!\left(
  \frac{A|_C\,\pi\,e|_C}{k_y|_C\,z^2}\,dz\right),
 \qquad
 \Res_0\!\left(
  \frac{A(0,z)\,\widetilde\pi}{b(z)^2}\,dz\right).
\end{equation}
The restrictions of $A$ to the two curves agree modulo $z^4$.
It follows from \eqref{trace:contact-expansion} that equality of their
residue functionals on all local test jets is exactly
\begin{equation}\label{trace:full-pp}
 \operatorname{pp}\!\left(\frac{\pi}{K(z)z^2}\right)
  =\operatorname{pp}\!\left(\frac{\widetilde\pi}{b(z)^2}\right),
 \qquad
 \widetilde\pi\equiv\frac{(b(z)/z)^2}{K(z)}\pi\pmod{z^2}.
\end{equation}
To display both coefficients, put
\[
 b=b_1z+b_2z^2+O(z^3),\quad K=K_0+K_1z+O(z^2),\quad
 \pi=\pi_0+\pi_1z+O(z^2).
\]
The respective principal parts are
\[
 \frac{\pi_0}{K_0}z^{-2}
 +\left(\frac{\pi_1}{K_0}-\frac{K_1\pi_0}{K_0^2}\right)z^{-1},
 \qquad
 \frac{\widetilde\pi_0}{b_1^2}z^{-2}
 +\left(\frac{\widetilde\pi_1}{b_1^2}
        -\frac{2b_2\widetilde\pi_0}{b_1^3}\right)z^{-1}.
\]
Hence
\begin{equation}\label{trace:two-coefficients}
 \begin{split}
 \widetilde\pi_0&=\frac{b_1^2}{K_0}\pi_0,\\
 \widetilde\pi_1&=\frac{b_1^2}{K_0}\pi_1+
       \left(\frac{2b_1b_2}{K_0}
                    -\frac{b_1^2K_1}{K_0^2}\right)\pi_0.
 \end{split}
\end{equation}
This invertible triangular transformation is the local expression of
the intrinsic residue-dual identification above; therefore its expressions
in different coordinates and frames agree.  In particular, no choice
of independent pointwise scalars is being substituted for a global map.

It remains to identify its metric.  In these frames write
$\eta_C=\mathcal E_C\,dz\,\mathfrak f^2$ and
$\eta_\Gamma=\mathcal E_\Gamma\,dz\,\mathfrak f^2$.
Tangency gives $\mathcal E_C(0)=\mathcal E_\Gamma(0)$.
Let $\ell$ and $\widetilde\ell$ be coefficient evaluations in $\tau$
and $\widetilde\tau$, respectively.  By
\eqref{trace:two-coefficients},
\[
 \iota^*\widetilde\ell=\alpha\ell,
 \qquad \alpha=\frac{b_1^2}{K_0}.
\]
Equations~\eqref{trace:dual-evaluation} and
\eqref{trace:frame-square} therefore give
\begin{equation}\label{trace:isometry-scalar}
 (\iota^*\widetilde\ell)^2
     =\alpha^2\frac{\mathcal E_C(0)}{s(0)}
     =-\mathcal E_\Gamma(0).
\end{equation}
The compatible bundle identifications give
$\widetilde\tau^2=\mathfrak f^2$.  Thus the final quantity in
\eqref{trace:isometry-scalar} is exactly $Q(\widetilde\ell)$, since
the section defining $Q$ multiplies to $-\eta_\Gamma/\omega$.
The evaluations at $Z$ have three distinct projective directions.
The difference of the two binary quadratic forms vanishes on those
three directions, and hence vanishes identically.  This proves the
isometry, with its sign and scalar fixed.
\end{proof}

\subsection{Meromorphic sections with the required residues}

The isometry now turns the six chosen vectors into prescribed pole
coefficients on $C$.  A global residue calculation will show that these
poles force the required local determinant normalization; no further
choice of regular parts will be needed.
We pass from the ambient coordinates used for adjunction to coordinates
adapted to $4L\simeq K_C$.  On a small disk centered at $q_j$ choose
a parameter $z_j$ and a frame $\ell_j^{\mathrm{fr}}$ of $L$ whose fourth
power maps to $dz_j$.  Such a frame exists by taking a holomorphic fourth
root of a unit.  Set
\[
 a_j=\xi/z_j,\qquad
 \mathfrak f_j=\ell_j^{\mathrm{fr}}/a_j,
 \qquad \tau_j=(\ell_j^{\mathrm{fr}})^2\mathfrak f_j.
\]
These are frames of $\OO_C(Z)$, $H$, and $\Lambda$, respectively.
If $r=\xi u$, its coefficient row in $\ell_j^{\mathrm{fr}}$ satisfies
$r=z_j u$, where here $u$ denotes its coefficient row in
$\mathfrak f_j$.  The fiber-linear function $t_j$ on $L^{-1}$
corresponding to $\ell_j^{\mathrm{fr}}$ gives compatible split
coordinates $(t_j,z_j)$.  Let $l_j\in W^*$ be coefficient evaluation
in $\tau_j$.  Since $\tau_j^2=dz_j\,\mathfrak f_j^2$,
Lemma~\ref{trace:pairing} reads in these coordinates as
\begin{equation}\label{trace:compatible-norm}
 l_j^2=(u_1u_2'-u_2u_1')(0).
\end{equation}
Primes in the rest of this section refer to $z_j$-derivatives on the
disk under consideration.

\begin{proposition}[Residue data]\label{prop:residue-data}
For the curve, line bundles, fixed fourth-root isomorphism, and six
vectors $w_j\in H^0(\Gamma,\Phi)^*$ of
Proposition~\ref{prop:curve-data}, one can choose the basis $u$ with
unchanged determinant so that $(u_1)_0$ is an unramified fiber of $f$
disjoint from $Z$.  Set $r=\xi u$ and define
\[
 v_j^*=\iota^*w_j\in W^*,\qquad V=\sum_{j=1}^6v_j^*.
\]
For any mutually disjoint compatible coordinate disks as above there
are nonzero constants $c_j$ and a global pair
\[
 P=(P_1,P_2)\in H^0(C,\theta+Z)^{\oplus2}
\]
with the following properties.  First,
\begin{equation}\label{trace:vector-identities}
 c_jl_j=v_j^*,\qquad
 V^2=90=9\deg L,\qquad (v_j^*)^2+v_j^*\cdot V=1.
\end{equation}
On disk $j$, using coefficient rows in the indicated frames, put
\begin{equation}\label{trace:local-data}
 m_j=\frac{c_j}{z_j},\qquad
 R_j=m_jr,\qquad
 W_j=R_j(0)\wedge R_j'(0),\qquad
 p_j=P-m_jr'.
\end{equation}
The pair $P$ has principal part $R_j(0)/z_j$ in the frame
$(\ell_j^{\mathrm{fr}})^2$ of $\theta$; the rows $R_j,p_j$ are
holomorphic on the disk, and $R_j(0)\ne0$.  Moreover,
\begin{equation}\label{trace:residue-identities}
 W_j=(v_j^*)^2,\qquad
 p_j(0)\wedge R_j(0)=1,\qquad
 \sum_{j=1}^6(W_j-1)=-9\deg L.
\end{equation}
Globally,
\begin{equation}\label{trace:J}
 J=u\wedge P\in H^0(C,\Lambda),\qquad
 \langle[J],[T]\rangle=-V([T])
       \quad\text{for every }T\in H^0(C,\Lambda).
\end{equation}
Changing $P$ by a holomorphic pair of sections of $\theta$ does not
change $[J]$.  The pole prescription is consistent under changes of
compatible coordinates.
\end{proposition}

\begin{proof}
A determinant-one change of basis preserves $\eta_C$ and
$\eta_\Gamma$, and hence all the forms already fixed.  Avoiding the
finitely many branch values and the image of $Z$, we may arrange that
the zero divisor of $u_1$ is the required fiber.  The isometry of
Lemma~\ref{trace:principal-parts} transports the equations for $w_j$
to \eqref{trace:vector-identities}.  It also transports their evaluation
directions.  Since each $l_j$ is nonzero, there is a unique nonzero
$c_j$ with $c_jl_j=v_j^*$.
Locally $R_j=c_ju$, so it is holomorphic and nonzero at the center, and
\eqref{trace:compatible-norm} gives
\[
 R_j(0)\wedge R_j'(0)=c_j^2l_j^2=(v_j^*)^2.
\]

Every prescribed set of simple principal parts of $\theta$ at $Z$
is realizable by a section of $\theta+Z$.  Indeed, the obstruction in
$H^1(C,\theta)$ pairs, by Serre duality, with
$H^0(C,K_C-\theta)=H^0(C,\theta)$ by the sum of residues.
All these test sections vanish at $Z$, by
Proposition~\ref{prop:curve-data}.  Their residue against any simple
principal part is therefore zero.  Apply this separately to the two
entries to obtain $P$ with the stated pole parts.  Since
$r=z_ju$, the row $m_jr'$ has exactly the same simple pole part as $P$;
thus $p_j$ is holomorphic.

The simple pole of $P$ at $q_j$ is proportional to $u(q_j)$.  Its
wedge with $u$ has zero residue coefficient, so $J=u\wedge P$ is
holomorphic there; away from $Z$ it is already holomorphic.  Hence it
is a global section of $\Lambda$.  To compute its class, fix
$T\in H^0(C,\Lambda)$ and consider the global meromorphic canonical
differential
\[
 \nu=\frac{P_1T}{u_1},
 \qquad \theta+(\theta+H)-H=2\theta=K_C.
\]
Its only possible poles lie in $Z$ and in the simple zero fiber of
$u_1$.  At $q_j$, the compatible frames make the canonical product
exactly multiplication of coefficients followed by $dz_j$, and therefore
\[
 \Res_{q_j}\nu=c_jl_j(T).
\]
At a zero $q$ of $u_1$, the section $u_2$ is nonzero and
\[
 J(q)=-u_2(q)P_1(q),\qquad
 \eta_C(q)=-u_2(q)\,du_1(q).
\]
Both expressions have the same minus sign.  In a local parameter,
the residue is consequently
\[
 \Res_q\nu=\left(\frac{JT}{\eta_C}\right)(q).
\]
The zero fiber is unramified, so the sum of these four residues is
the value of $\Tr(JT/\eta_C)$ there.  This trace is a global
holomorphic function on $\PP^1$ by Lemma~\ref{trace:pairing}, hence
is constant.  The residue theorem gives
\[
 \langle[J],[T]\rangle
       =-\sum_jc_jl_j(T)=-V([T]),
\]
which proves \eqref{trace:J} with the stated sign.

On a fixed disk suppress the subscript $j$.  Since $r=zu$,
\[
 P=p+\frac{c}{z}u+cu',\qquad
 cJ=R\wedge p+R\wedge R'.
\]
In particular,
\[
 c_jl_j(J)=R_j(0)\wedge p_j(0)+W_j.
\]
On the other hand, the metric identification in \eqref{trace:J} gives
\[
 c_jl_j(J)=-v_j^*\cdot V=W_j-1.
\]
Thus $R_j(0)\wedge p_j(0)=-1$, or equivalently
$p_j(0)\wedge R_j(0)=1$.  Summing
$W_j-1=-v_j^*\cdot V$ proves the last identity in
\eqref{trace:residue-identities}.  Notice also the useful local expansion
\begin{equation}\label{trace:P-constant}
 P=\frac{R_j(0)}{z_j}
       +p_j(0)+2R_j'(0)+O(z_j).
\end{equation}
If $P$ is changed by a holomorphic pair $S$, the equality
$H^0(C,\theta)=\xi^2H^0(C,2H)$ shows that
$u\wedge S\in\xi^2H^0(C,3H)$; thus $[J]$ is unchanged.

For completeness, let $\zeta=\psi(z)$ be another centered parameter,
and let the corresponding compatible frame of $L$ be $g(z)\ell^{\mathrm{fr}}$.
Compatibility says $g^4=\psi'$.  The associated changes of the other
frames are
\[
 a_{\rm new}=\frac{z}{\psi(z)}a,\qquad
 \mathfrak f_{\rm new}=g\frac{\psi(z)}{z}\mathfrak f,
 \qquad \tau_{\rm new}=g^3\frac{\psi(z)}{z}\tau.
\]
Writing $g_0=g(0)$, it follows that
$l_{\rm new}=g_0^{-7}l$, so the same vector $v^*$ prescribes
$c_{\rm new}=g_0^7c$.  Also
\[
 r_{\rm new}=g^{-1}r,\qquad
 R_{\rm new}(0)=g_0^2R(0).
\]
The coefficient row of the global $\theta$-section transforms as
$P_{\rm new}=g^{-2}P$.  Its coefficient of $\zeta^{-1}$ is
$\psi'(0)g_0^{-2}R(0)=g_0^2R(0)$, exactly the newly prescribed
$R_{\rm new}(0)$.  Thus the same global $P$ realizes the new
prescription.  The row $R_{\rm new}$ is a scalar multiple of $R$,
with scalar $g_0^2$ at the center; differentiating with respect to
$\zeta$ shows $W_{\rm new}=g_0^4W/\psi'(0)=W$.
Finally, subtracting $m_{\rm new}\,d r_{\rm new}/d\zeta$ from
$P_{\rm new}$ shows that
$p_{\rm new}(0)-g_0^{-2}p(0)$ is proportional to $R(0)$.
Hence $p_{\rm new}(0)\wedge R_{\rm new}(0)=p(0)\wedge R(0)$.
All the local conclusions are therefore compatible with the fixed
fourth-root structure.
\end{proof}
